%\documentclass[10pt,respuestas,a4]{examenEPN}
\documentclass[10pt,a4]{examenEPN0}
\examen{Examen 01}
\usepackage{comandosEPN}

\usepackage{empheq}

\newcommand{\itemcp}{\stepcounter{enumi}\item[ \textcolor{red}{\theenumi.$^*$}] }

\newcommand{\iitemcp}{\stepcounter{enumii}\item[ \textcolor{red}{\theenumii)$^*$}] }


\materia{Álgebra Lineal}
%\autor{Departamento de Formación Básica}
\autor{GR \hspace{1cm}.}
%\fecha{Semestre 2026-A}
\fecha{Nombre: \hspace{9cm} Cédula:}
\logouno[1.6cm]{Logo_EPN.eps}
\logodos[2.6cm]{Logo_DFB.eps}

\newcommand{\la}{\lambda}

%Datos
\newcommand{\datosAlgebra}[2]{
\EnExamen
{
\vspace{5mm}
    \noindent
    Nombres: \hrulefill\
    Carrera: \rule{4.5cm}{0.4pt}\
    \ifthenelse{\equal{#1}{}}%
    {Paralelo: \underline{\parbox{1.8cm}{\parbox{1cm}{}}}\\[3mm]}
    {Paralelo: \underline{\parbox{1.8cm}{\centering #1}}\\[3mm]}
    C\'edula: \rule{5cm}{0.4pt}\
    Firma: \hrulefill\
    Fecha: \underline{\parbox{2.1cm}{\centering #2}}
}}


% \newcommand{\indicacionesPrueba}{%
% \EnExamen{
% \begin{indicaciones}
%     \item La prueba tiene una duraci\'on  de \textbf{120 minutos.}
%     \item No se permite la utilizaci\'on de dispositivos electr\'onicos (celular, tablets, calculadora, etc.)
%     \item Leer detenidamente los enunciados, \textbf{justificar la respuesta para acceder al puntaje total}.
%     \item La prueba es completamente individual. Los estudiantes deberán dejar sus teléfonos celulares guardados y no podrán tenerlos consigo (por ejemplo, en los bolsillos de su ropa) durante toda la evaluación. En caso de detectarse algún acto de deshonestidad académica o el incumplimiento de esta norma, no podrán seguir rindiendo la prueba y su calificación será \textbf{cero}.
% \end{indicaciones}
% \pagestyle{empty}
% }}

\newcommand{\indicacionesPrueba}{%
\EnExamen{
\begin{indicaciones}
    \item La prueba tiene una duraci\'on  de \textbf{60 minutos.}
    \item Leer detenidamente los enunciados, \textbf{justificar la respuesta} para acceder al puntaje total.
    \item No se permite la utilizaci\'on de dispositivos electr\'onicos (celular, tablets, calculadora, etc.)
    \item La prueba es completamente individual
    \item No se permite portar teléfonos celulares, relojes inteligentes ni otros dispositivos electrónicos durante el examen, aun cuando permanezcan en los bolsillos. Estos dispositivos deberán guardarse en la mochila antes del inicio de la evaluación y esta deberá mantenerse cerrada durante todo el examen.
    \item He leído las indicaciones y acepto que el incumplimiento de estas se penaliza con \textbf{cero}. Firma:...................................

    
\end{indicaciones}
\pagestyle{empty}
}}

\newcommand{\indicacionesExamen}{
\EnExamen{\begin{indicaciones}
    \item La evaluación tiene una duraci\'on  de \textbf{120 minutos.}
    \item No se permite la utilizaci\'on de dispositivos electr\'onicos (celular, tablets, calculadora, etc.)
    %\item Leer detenidamente los enunciados, justificar la respuesta para acceder al puntaje total.
    \item El examen es completamente individual. En caso de detectarse alg\'un acto de deshonestidad acad\'emica no podr\'a seguir rindiendo la prueba y su calificaci\'on ser\'a \textbf{cero}.
\end{indicaciones}
}\pagestyle{empty}
}

%% Imprimir las una hoja A4, dos veces en A5
\newcommand{\imprimirDoble}
{%
\usepackage{pgfmorepages}
\pgfmorepagesloadextralayouts
\pgfpagesuselayout{repeated 2-up}[a4paper,landscape]
\thispagestyle{empty}
}


%Este sí se puede comentar
\newenvironment{solucion}{\color{blue}\proof[Solución]}{\endproof\color{black}}




% \renewcommand{\@autor}{
% ---\\
% \phantom{.} 
% ...}
% \makeatother

%\imprimirDoble
\begin{document}
\encabezado
\datosAlgebra{GR5E}{19/11/2026}
%Poner el grupo y la fecha, o dejar las llaves vacías.
\indicacionesPrueba
\begin{preguntas}


%%%%%%%%%%%%%%%%%%%%%%%%%%%%%%%%%

%%%%%%%%%%%%%%%%%%%%%%%%%%%%%%%%%%%%%%%%%
\item Una empresa alquila proyectores, pantallas y sistemas de sonido con tarifas diarias constantes. Alquilar una unidad de cada equipo durante un día cuesta \$65. Dos proyectores y una pantalla, durante el mismo tiempo, cuestan \$80. La tarifa diaria de una pantalla supera en \$5 la de un sistema de sonido. Determine las tres tarifas.

\begin{enumerate}
    \item Defina las tres variables, plantee el sistema de ecuaciones lineales que modela la situación y escriba su representación matricial $A\mathbf{x}=b$, identificando $A$, $\mathbf{x}$ y $b$. \puntaje{2.0}
    \item Resuelva el sistema e indique las tarifas detalladas en el problema. \puntaje{2.0}
\end{enumerate}
\newpage

%%%%%%%%%%%%%%%%%%%%%%%%%%%%%%%%%
\small
\item Dado el sistema de ecuaciones

$$
    \left\{
    \begin{alignedat}{4}
        x_1  & + \alpha x_2       & + 2x_3              &= 1 \\
        -x_1 & - 2x_2             & + (\alpha -2)x_3    &= -1 \\
             & (\alpha - 2) x_2   & + (2\alpha +2)x_3  &= 1
    \end{alignedat}
    \right.
$$
\begin{enumerate}
    \item Calcule el determinante de la matriz de coeficientes $A$ y determine los valores de $\alpha\in\mathbb{R}$ para los cuales el sistema de ecuaciones tiene solución única. \puntaje{2.0}
    \item Para $\alpha=2$, halle la solución del sistema. \puntaje{2.0} 
\end{enumerate}
\newpage


%%%%%%%%%%%%%%%%%%%%%%%%%%%%%%%%%%%%%%%
\item Considere S y T subespacios vectoriales de $\R^{2\times2}$,
\begin{center}
$S=\{A : a_{12}=a_{21}\},\qquad
T=\{A : a_{11}+a_{22}=0,\ a_{11}-a_{12}=0\}.$
\end{center}
\begin{enumerate}
\item Hallar $S\cap T$ de forma explícita. \puntaje{2.0}
\item Hallar $k\in\R$ tal que $A=\begin{pmatrix}k+1&3\\2k-1&k-5\end{pmatrix}\in S\cap T$. \puntaje{2.0}
\end{enumerate}
\newpage
%%%%%%%%%%%%%%%%%%%%%%%%%%%%%%%%%%%%%%%%%
\item Considere los vectores de \(\mathbb{R}^3\):
\[
u = (1, 2, -1), \quad v = (2, 1, 3), \quad w = (1, -1, m),
\]
con \(m \in \mathbb{R}\).

\begin{enumerate}
    \item Calcule el \(\operatorname{span}\{u, v\}\) e interprete geométricamente si el espacio generado es una recta o un plano en \(\mathbb{R}^3\). Justifique. \puntaje{2.0} 
    \item Determine \textbf{para qué valor(es) de \(m\)} el vector \(w\) pertenece a \(\operatorname{span}\{u, v\}\). Exprese $w$ como combinación lineal de $u$ y $v$. \puntaje{1.0} 
    \item Dado el vector \(x = (-1, 2, 2)\). Determine si el vector \(x\) pertenece a \(\operatorname{span}\{u, v\}\). \puntaje{1.0} 
   
\end{enumerate}
\newpage

%%%%%%%%%%%%%%%%%%%%%%%%%%%%%%%%%%%%%%%
\item (Ejercicio del deber) Sea $V = \mathcal{P}_3(\R)$ el espacio de polinomios de grado menor o igual a tres con coeficientes reales. Determine si el siguiente conjunto es un subespacio de $V$: \puntaje{4.0} 
    \[
    W = \left\{ p(x) \in \mathcal{P}_3(\R) : p(2) = 0 \quad \text{y} \quad p'(0) = 0 \right\}
    \]      




\end{preguntas}

\end{document}